\section{Conclusion}
\label{sec:conclusion}
In this work, we presented \textbf{V-Dreamer}, a fully automated framework for synthesizing large-scale, physically grounded simulation environments and executable manipulation trajectories for generalist robot learning. By integrating Large Language Models with advanced 2D/3D and video generative models, V-Dreamer translates open-vocabulary instructions into diverse interactive scenes and coherent manipulation behaviors. Unlike prior generative pipelines that primarily focus on scene synthesis, V-Dreamer tackles the full generation loop, producing training-ready data that is directly usable for policy learning. 

Extensive experiments show that imitation policies trained exclusively on V-Dreamer-generated data achieve robust performance and zero-shot generalization to geometrically unseen objects in simulation. Furthermore, with photo-conditioned scene alignment, these policies transfer zero-shot to real hardware. These results suggest that the trajectories generated by V-Dreamer are not merely visually plausible, but physically executable and directly useful for robot control, highlighting the promise of full-cycle generative pipelines in alleviating the data bottleneck in robotic manipulation.

\textbf{Limitations and Future Work.}
An important next step is to develop automated, physics-aware trajectory filtering to reduce the impact of low-quality generations. In addition, the current framework is limited to rigid-body tabletop manipulation. Extending it to articulated and deformable objects is a key direction for future work. Finally, incorporating tighter physical feedback during generation and grounding may further improve physical realism and sim-to-real robustness.